\subsection*{Conditioning is a controllability mechanism, not an extrapolation mechanism}

The experiments support a simple but consequential conclusion. Continuous conditioning is useful: it improves target adherence and preserves usable generation closer to the observed boundary than the fixed-threshold binary baseline. But the same mechanism does not create reliable candidates beyond the boundary. The zero frontier-crossing result at Q90, Q100 and both train-max offsets is unchanged by moving from one property to six.

This observation also clarifies how to compare the present benchmark with fragment-based or multi-property molecular generators such as MolMiner. In-support enrichment, target ordering and Pareto reweighting are valuable capabilities, but none is by itself evidence of numerical OOD extrapolation. A fair comparison must freeze the training support, define the frontier before generation and audit candidate predictions against that boundary.

The practical implication is that the next gain is unlikely to come from conditioning design alone. The most promising directions are: (i) adding genuinely labelled conductivity data from harmonized experimental sources, rather than using an unlabelled DFT database as if it supplied conductivity; (ii) testing fragment-based or otherwise richer representations; and (iii) coupling generation to active-learning or MD validation so that new frontier regions become labelled and feed back into the model.

\placeholder{Add literature citations and a quantitative comparison to the final MolMiner manuscript once the bibliography and exact comparator protocol are frozen.}
